
\section{Advanced Activities}

The following activities are intended for strong master's or PhD students.
They integrate the main ideas developed throughout the course through
active-suspension and wind--wave-energy applications. The activities emphasize
fair architecture comparison, robustness, surrogate and multi-fidelity
modeling, implementable hardware, and reproducible validation. MATLAB, Python,
OpenMDAO, GPOPS-II, Dymos, IPOPT, SNOPT, or equivalent tools may be used where
specified.

\subsection*{Problem 1: Full Active-Suspension CCD Benchmark}

Consider the quarter-car suspension model
\begin{align}
m_s\ddot{z}_s
&=
-k_s(z_s-z_u)
-c_s(\dot{z}_s-\dot{z}_u)
+f_a,
\\
m_u\ddot{z}_u
&=
k_s(z_s-z_u)
+c_s(\dot{z}_s-\dot{z}_u)
-k_t(z_u-z_r)
-f_a,
\end{align}
with actuator dynamics
\begin{equation}
\dot{f}_a
=
\omega_a
\left(
f_c-f_a
\right).
\end{equation}

Use
\begin{equation}
m_s=300~\mathrm{kg},
\qquad
m_u=40~\mathrm{kg},
\qquad
k_t=190\,000~\mathrm{N/m},
\end{equation}
and
\begin{equation}
\omega_a=60~\mathrm{rad/s}.
\end{equation}

The plant-design variables are
\begin{equation}
10\,000\leq k_s\leq40\,000~\mathrm{N/m},
\end{equation}
\begin{equation}
500\leq c_s\leq3500~\mathrm{N\,s/m},
\end{equation}
and
\begin{equation}
500\leq F_{\max}\leq4000~\mathrm{N}.
\end{equation}

Use the causal feedback controller
\begin{equation}
f_c(t)
=
\operatorname{sat}
\left[
-K_1(z_s-z_u)
-K_2(\dot{z}_s-\dot{z}_u)
-K_3(z_u-z_r)
-K_4\dot{z}_u,
F_{\max}
\right].
\end{equation}

The controller gains satisfy
\begin{equation}
0\leq K_1\leq60\,000,
\qquad
0\leq K_2\leq6000,
\end{equation}
\begin{equation}
0\leq K_3\leq80\,000,
\qquad
0\leq K_4\leq6000.
\end{equation}

The road input is a half-cosine bump:
\begin{equation}
z_r(t)
=
\begin{cases}
\dfrac{h_b}{2}
\left[
1-\cos
\left(
\dfrac{2\pi(t-t_b)}{T_b}
\right)
\right],
&
t_b\leq t\leq t_b+T_b,
\\[1em]
0,
&
\text{otherwise},
\end{cases}
\end{equation}
with
\begin{equation}
h_b=0.05~\mathrm{m},
\qquad
t_b=0.25~\mathrm{s},
\qquad
T_b=0.10~\mathrm{s}.
\end{equation}

Use the system-level objective
\begin{align}
J
=&
\frac{1}{T}
\int_0^T
\Bigg[
0.45
\left(
\frac{\ddot{z}_s}{2.5}
\right)^2
+
0.20
\left(
\frac{z_u-z_r}{0.015}
\right)^2
\nonumber\\
&\qquad
+
0.20
\left(
\frac{z_s-z_u}{0.05}
\right)^2
+
0.10
\left(
\frac{f_c}{3000}
\right)^2
\Bigg]dt
\nonumber\\
&+
0.05
\left(
\frac{F_{\max}}{3000}
\right)^2
+
0.02
\left(
\frac{k_s}{25\,000}
\right)^2
+
0.02
\left(
\frac{c_s}{2000}
\right)^2,
\end{align}
with
\begin{equation}
T=2~\mathrm{s}.
\end{equation}

Impose
\begin{equation}
|z_s-z_u|\leq0.08~\mathrm{m},
\end{equation}
\begin{equation}
|z_u-z_r|\leq0.025~\mathrm{m},
\end{equation}
\begin{equation}
|f_c|\leq F_{\max},
\qquad
|f_a|\leq F_{\max}.
\end{equation}

\begin{enumerate}
    \item Derive the five-state nonlinear closed-loop model.

    \item Formulate the single-pass sequential design:
    \begin{enumerate}
        \item optimize $k_s$, $c_s$, and $F_{\max}$ using a fixed nominal
        controller;
        \item freeze the plant and optimize $K_1,\ldots,K_4$.
    \end{enumerate}

    \item Formulate the nested design
    \begin{equation}
    \min_{k_s,c_s,F_{\max}}
    \left[
    \min_{K_1,\ldots,K_4}
    J
    \right].
    \end{equation}

    \item Formulate the simultaneous CCD problem.

    \item Solve all three formulations using identical dynamics, constraints,
    integration tolerances, initial conditions, and objective weights.

    \item Use at least ten initial guesses for the nested and simultaneous
    formulations.

    \item Compare the best designs in terms of:
    \begin{equation}
    J^*,
    \quad
    k_s^*,
    \quad
    c_s^*,
    \quad
    F_{\max}^*,
    \quad
    K_1^*,\ldots,K_4^*.
    \end{equation}

    \item Verify every candidate design using an independent high-accuracy
    forward simulation and a time grid at least ten times finer than the
    optimization grid.

    \item Explain any disagreement between the nested and simultaneous
    results in terms of local minima, inner-loop convergence, derivatives,
    scaling, and solver termination criteria.
\end{enumerate}


\subsection*{Problem 2: Robust Suspension Design using Expected Value, Worst Case, and CVaR}

Use the suspension model from Problem 1. Define the uncertain sprung mass
\begin{equation}
m_s
\in
\{240,\ 300,\ 360\}~\mathrm{kg},
\end{equation}
and the uncertain bump height
\begin{equation}
h_b
\in
\{0.03,\ 0.05,\ 0.07\}~\mathrm{m}.
\end{equation}

This produces nine equally probable scenarios. Let the scenario objective be
\begin{equation}
J_s(\mathbf{x}_p,\mathbf{x}_c),
\qquad
s=1,\ldots,9.
\end{equation}

\begin{enumerate}
    \item Formulate the expected-value problem
    \begin{equation}
    \min_{\mathbf{x}_p,\mathbf{x}_c}
    \frac{1}{9}
    \sum_{s=1}^{9}J_s.
    \end{equation}

    \item Formulate the worst-case problem using an epigraph variable $\eta$:
    \begin{align}
    \min_{\mathbf{x}_p,\mathbf{x}_c,\eta}
    \quad
    & \eta,
    \\
    \text{subject to}
    \quad
    & J_s\leq\eta,
    \qquad
    s=1,\ldots,9.
    \end{align}

    \item Formulate the combined mean--worst-case problem
    \begin{equation}
    \min
    \quad
    \frac{1}{9}
    \sum_{s=1}^{9}J_s
    +
    \lambda\eta,
    \end{equation}
    for
    \begin{equation}
    \lambda
    \in
    \{0,\ 0.1,\ 0.5,\ 1,\ 2\}.
    \end{equation}

    \item Formulate the empirical CVaR objective at confidence level
    \begin{equation}
    \alpha=0.9
    \end{equation}
    using auxiliary variables $\xi_s$:
    \begin{equation}
    \operatorname{CVaR}_{\alpha}
    =
    \eta
    +
    \frac{1}{(1-\alpha)N_s}
    \sum_{s=1}^{N_s}\xi_s,
    \end{equation}
    subject to
    \begin{equation}
    \xi_s\geq J_s-\eta,
    \qquad
    \xi_s\geq0.
    \end{equation}

    \item Require every path constraint to hold in all nine scenarios.

    \item Solve the nominal, expected-value, worst-case, and CVaR designs.

    \item Compare:
    \begin{equation}
    J_{\mathrm{nominal}},
    \qquad
    \mathbb{E}[J],
    \qquad
    \max_sJ_s,
    \qquad
    \operatorname{CVaR}_{0.9}.
    \end{equation}

    \item Quantify the nominal-performance penalty of robustness.

    \item Determine which plant and controller variables change most strongly
    when moving from nominal to robust design.

    \item Explain why a robust design can have a worse nominal objective but
    still be the more credible engineering design.
\end{enumerate}


\subsection*{Problem 3: Surrogate-Assisted Wind-Turbine CCD with Adaptive Infill}

Consider a simplified floating-wind CCD problem with plant variables
\begin{equation}
\mathbf{x}_p
=
\begin{bmatrix}
D_p&
K_m&
C_m
\end{bmatrix}^{T},
\end{equation}
where $D_p$ is a platform-dimension parameter, $K_m$ is an effective mooring
stiffness, and $C_m$ is an effective mooring damping parameter.

The controller variables are
\begin{equation}
\mathbf{x}_c
=
\begin{bmatrix}
K_{\Omega}&
K_{\theta}&
K_{\dot{\theta}}
\end{bmatrix}^{T}.
\end{equation}

A high-fidelity simulator returns
\begin{equation}
\mathbf{y}
=
\begin{bmatrix}
P_{\mathrm{mean}}&
\theta_{\mathrm{rms}}&
M_{\mathrm{tower,max}}&
T_{\mathrm{moor,max}}
\end{bmatrix}^{T}.
\end{equation}

The system-level objective is
\begin{align}
J_{\mathrm{HF}}
=&
-
\frac{P_{\mathrm{mean}}}{P_{\mathrm{ref}}}
+
0.3
\left(
\frac{\theta_{\mathrm{rms}}}{\theta_{\mathrm{ref}}}
\right)^2
\nonumber\\
&+
0.2
\left(
\frac{M_{\mathrm{tower,max}}}{M_{\mathrm{ref}}}
\right)^2
+
0.1
\left(
\frac{T_{\mathrm{moor,max}}}{T_{\mathrm{ref}}}
\right)^2
+
C_{\mathrm{plant}}(\mathbf{x}_p).
\end{align}

\begin{enumerate}
    \item Construct a Latin-hypercube design of experiments containing at
    least 150 plant--controller samples.

    \item Train separate Gaussian-process or radial-basis surrogates for all
    four outputs.

    \item Split the data into training, validation, and test sets, and report:
    \begin{equation}
    R^2,
    \qquad
    \mathrm{RMSE},
    \qquad
    \max|e|.
    \end{equation}

    \item Formulate the surrogate CCD problem, including high-fidelity output
    constraints
    \begin{equation}
    M_{\mathrm{tower,max}}\leq M_{\mathrm{allow}},
    \end{equation}
    and
    \begin{equation}
    T_{\mathrm{moor,max}}\leq T_{\mathrm{allow}}.
    \end{equation}

    \item Solve the surrogate optimization from at least twenty initial
    guesses.

    \item Evaluate the best five surrogate candidates in the high-fidelity
    model.

    \item Define a combined infill criterion
    \begin{equation}
    \mathcal{I}(\mathbf{z})
    =
    \widehat{J}(\mathbf{z})
    -
    \kappa\sigma_J(\mathbf{z})
    +
    \rho\,V(\mathbf{z}),
    \end{equation}
    where $\sigma_J$ is predictive uncertainty and $V$ is a predicted
    constraint-violation penalty.

    \item Add at least five adaptive infill points per iteration and repeat
    until the high-fidelity objective changes by less than
    \begin{equation}
    0.5\%.
    \end{equation}

    \item Compare the surrogate optimum and final high-fidelity optimum.

    \item Explain why a surrogate that predicts only annual energy or mean
    power is insufficient when transient loads appear in the constraints.
\end{enumerate}


\subsection*{Problem 4: Multi-Fidelity CCD with a Trust-Region Correction Model}

Let
\begin{equation}
J_L(\mathbf{z})
\end{equation}
be a low-fidelity CCD objective and
\begin{equation}
J_H(\mathbf{z})
\end{equation}
the corresponding high-fidelity objective, where
\begin{equation}
\mathbf{z}
=
\begin{bmatrix}
\mathbf{x}_p\\
\mathbf{x}_c
\end{bmatrix}.
\end{equation}

Define the discrepancy
\begin{equation}
\delta(\mathbf{z})
=
J_H(\mathbf{z})-J_L(\mathbf{z}).
\end{equation}

At iteration $k$, construct the corrected model
\begin{equation}
m_k(\mathbf{z})
=
J_L(\mathbf{z})
+
\widehat{\delta}_k(\mathbf{z}),
\end{equation}
and solve
\begin{equation}
\min_{\mathbf{z}}
\quad
m_k(\mathbf{z})
\end{equation}
subject to
\begin{equation}
\left\|
\mathbf{z}-\mathbf{z}^{(k)}
\right\|_2
\leq
\Delta_k.
\end{equation}

\begin{enumerate}
    \item Derive the predicted reduction
    \begin{equation}
    \Delta m_k
    =
    m_k(\mathbf{z}^{(k)})
    -
    m_k(\mathbf{z}^{(k)}+\mathbf{s}_k).
    \end{equation}

    \item Derive the actual high-fidelity reduction
    \begin{equation}
    \Delta J_{H,k}
    =
    J_H(\mathbf{z}^{(k)})
    -
    J_H(\mathbf{z}^{(k)}+\mathbf{s}_k).
    \end{equation}

    \item Define the trust-region ratio
    \begin{equation}
    \rho_k
    =
    \frac{\Delta J_{H,k}}
    {\Delta m_k}.
    \end{equation}

    \item Propose explicit rules for:
    \begin{enumerate}
        \item accepting or rejecting the trial design;
        \item shrinking the trust region;
        \item expanding the trust region.
    \end{enumerate}

    \item Implement the method for a low-fidelity suspension or wind model
    and a higher-fidelity validation model.

    \item Use both additive correction
    \begin{equation}
    J_H\approx J_L+\widehat{\delta}
    \end{equation}
    and multiplicative correction
    \begin{equation}
    J_H\approx\widehat{\beta}J_L.
    \end{equation}

    \item Compare the two correction strategies in terms of:
    \begin{equation}
    \text{high-fidelity calls},
    \quad
    \text{final objective},
    \quad
    \text{constraint satisfaction}.
    \end{equation}

    \item Determine whether the low-fidelity optimum lies inside a region
    where the correction model is accurate.

    \item Explain why unconstrained optimization of a corrected surrogate can
    be unreliable without a trust region or another mechanism that controls
    extrapolation.
\end{enumerate}


\subsection*{Problem 5: From Optimized Suspension to Implementable Hardware}

Use the optimized active-suspension design from Problem 1. Replace the ideal
implementation with the sampled-data controller
\begin{equation}
f_{c,k}
=
\operatorname{sat}
\left[
-K\widehat{\mathbf{x}}_k,
F_{\max}
\right],
\end{equation}
where the controller is updated every
\begin{equation}
T_s
\end{equation}
seconds and held constant between updates.

The sensor measurements satisfy
\begin{equation}
\mathbf{y}_k
=
C\mathbf{x}_k
+
\mathbf{v}_k,
\end{equation}
where $\mathbf{v}_k$ is zero-mean measurement noise.

The command is delayed by
\begin{equation}
n_d
\end{equation}
sampling intervals, and the actuator also satisfies
\begin{equation}
|\dot{f}_a|
\leq
R_{\max}.
\end{equation}

\begin{enumerate}
    \item Derive a discrete-time model using zero-order hold.

    \item Construct a state observer or Kalman filter for the available
    measurements.

    \item Implement the delayed command
    \begin{equation}
    f_{c,k}^{\mathrm{applied}}
    =
    f_{c,k-n_d}.
    \end{equation}

    \item Impose the actuator-rate limit using
    \begin{equation}
    |f_{a,k+1}-f_{a,k}|
    \leq
    R_{\max}T_s.
    \end{equation}

    \item Evaluate
    \begin{equation}
    T_s
    \in
    \{0.001,\ 0.005,\ 0.01,\ 0.02,\ 0.05\}~\mathrm{s}
    \end{equation}
    and
    \begin{equation}
    n_d
    \in
    \{0,\ 1,\ 2,\ 5\}.
    \end{equation}

    \item Determine the largest sampling time and delay for which all path
    constraints remain satisfied.

    \item Add a passive-safety requirement: when the actuator is disabled,
    the passive suspension must satisfy
    \begin{equation}
    |z_s-z_u|
    \leq
    0.10~\mathrm{m}.
    \end{equation}

    \item Re-optimize the plant and controller while including sampling time,
    command delay, observer dynamics, force-rate limits, and passive safety.

    \item Compare the ideal and implementable CCD designs in terms of:
    \begin{equation}
    J,
    \quad
    k_s,
    \quad
    c_s,
    \quad
    F_{\max},
    \quad
    \text{constraint margin}.
    \end{equation}

    \item Design a software-in-the-loop, hardware-in-the-loop, and bench-test
    validation sequence for the final design.
\end{enumerate}


\subsection*{Problem 6: Independent Reproduction and Validation Audit}

A research team reports an optimized CCD design
\begin{equation}
\mathbf{z}^*
=
\begin{bmatrix}
\mathbf{x}_p^*\\
\mathbf{x}_c^*
\end{bmatrix},
\end{equation}
with reported objective
\begin{equation}
J_{\mathrm{reported}},
\end{equation}
maximum path-constraint violation below
\begin{equation}
10^{-6},
\end{equation}
and a claimed
\begin{equation}
18\%
\end{equation}
improvement over a sequential baseline.

You are asked to perform an independent reproduction and validation audit.

\begin{enumerate}
    \item Define a reproducibility package containing:
    \begin{enumerate}
        \item governing equations;
        \item units and variable definitions;
        \item initial and boundary conditions;
        \item objective weights;
        \item all constraints and bounds;
        \item software versions;
        \item solver settings;
        \item initialization;
        \item random seeds;
        \item raw and processed data.
    \end{enumerate}

    \item Recompute the reported objective using an independent implementation.

    \item Define the relative objective discrepancy
    \begin{equation}
    e_J
    =
    \frac{
    |J_{\mathrm{reproduced}}-J_{\mathrm{reported}}|
    }{
    \max(1,|J_{\mathrm{reported}}|)
    }.
    \end{equation}

    \item Independently simulate the reported control law and compute the
    maximum trajectory discrepancy
    \begin{equation}
    e_x
    =
    \max_t
    \left\|
    \mathbf{x}_{\mathrm{reported}}(t)
    -
    \mathbf{x}_{\mathrm{independent}}(t)
    \right\|_{\infty}.
    \end{equation}

    \item Re-evaluate every path constraint on a dense grid and report
    \begin{equation}
    v_{\max}
    =
    \max_{t,j}
    g_j
    \left(
    \mathbf{x}(t),
    \mathbf{u}(t)
    \right).
    \end{equation}

    \item Repeat the optimization using:
    \begin{enumerate}
        \item at least ten initial guesses;
        \item two mesh densities;
        \item two NLP solvers or algorithms;
        \item two derivative methods.
    \end{enumerate}

    \item Reconstruct the sequential baseline using exactly the same model,
    objective, constraints, tolerances, and information assumptions as the
    CCD problem.

    \item Compute the verified improvement
    \begin{equation}
    \eta_{\mathrm{verified}}
    =
    \frac{
    J_{\mathrm{seq,verified}}
    -
    J_{\mathrm{CCD,verified}}
    }{
    J_{\mathrm{seq,verified}}
    }.
    \end{equation}

    \item Test both designs under uncertainty and at a higher model-fidelity
    level.

    \item Define explicit acceptance criteria for:
    \begin{equation}
    e_J,
    \qquad
    e_x,
    \qquad
    v_{\max},
    \qquad
    \eta_{\mathrm{verified}}.
    \end{equation}

    \item Classify the final conclusion as one of the following:
    \begin{enumerate}
        \item numerically reproduced and physically validated;
        \item numerically reproduced but not physically validated;
        \item partially reproduced;
        \item not reproduced.
    \end{enumerate}

    \item Explain why an optimizer success message and a low reported
    objective are insufficient evidence for a credible CCD result.
\end{enumerate}
